\subsection{\texorpdfstring{ALGEBRAS OF TYPE \(\mathbf{A}_l\) ( \(l \geq 1\) )}{ALGEBRAS OF TYPE \textbackslash mathbf\{A\}\_l ( l \textbackslash geq 1 )}}

\begin{enumerate}
\def\labelenumi{(\Roman{enumi})}
\tightlist
\item
  Let V be a vector space of dimension \(l+1\) over k, and let \(\mathfrak{g}\) be the algebra \(\mathfrak{sl}(V)\) of endomorphisms of V of trace zero. Let \((e_{i})_{1\leq i\leq l+1}\) be a basis of V; the map which associates to an element of \(\mathfrak{g}\) its matrix with respect to this basis is an identification of \(\mathfrak{g}\) with the algebra \(\mathfrak{sl}(l+1,k)\) of matrices of trace zero. We know that \(\mathfrak{g}\) is semi-simple (Chap. I, §6, no. 7, Prop. 8).
\end{enumerate}

Recall (Algebra, Chap. II, §10, no. 3) that \(E_{ij}\) denotes the matrix \((\alpha_{mp})\) such that \(\alpha_{ij} = 1\) and \(\alpha_{mp} = 0\) for \((m,p) \neq (i,j)\) . The matrices

\[
\begin{array}{l l} E _ {i j} & (1 \leq i, j \leq l + 1, i \neq j) \\ E _ {i, i} - E _ {i + 1, i + 1} & (1 \leq i \leq l) \end{array}
\]

form a basis of \(\mathfrak{g}\). Hence

\[
\dim \mathfrak {g} = l (l + 2).
\]

Let \(\hat{\mathfrak{h}}\) be the set of diagonal elements of \(\mathfrak{gl}(l + 1,k)\) ; the sequence \((E_{ii})_{1\leq i\leq l + 1}\) is a basis of the vector space \(\hat{\mathfrak{h}}\) ; let \((\hat{\varepsilon}_i)_{1\leq i\leq l + 1}\) be the basis of \(\hat{\mathfrak{h}}^*\) dual to \((E_{ii})_{1\leq i\leq l + 1}\) . For all \(h\in \hat{\mathfrak{h}}\) ,

\[
[ h, E _ {i j} ] = (\hat {\varepsilon} _ {i} (h) - \hat {\varepsilon} _ {j} (h)) E _ {i j}\tag{1}
\]

by Chap. I, §1, no. 2, formulas (5). Let \(\mathfrak{h}\) be the set of elements of \(\hat{\mathfrak{h}}\) of trace zero, and put \(\varepsilon_i = \hat{\varepsilon}_i|\mathfrak{h}\) . Then \(\mathfrak{h}\) is a Cartan subalgebra of \(\mathfrak{g}\) (Chap. VII, §2, no. 1, Example 4). Relation (1) proves that this Cartan subalgebra is splitting, and that the roots of \((\mathfrak{g},\mathfrak{h})\) are the \(\varepsilon_i - \varepsilon_j\) ( \(i\neq j\) ). Let \(\hat{\mathfrak{h}}_0^*\) be the set of elements of \(\hat{\mathfrak{h}}^*\) the sum of whose coordinates with respect to \((\hat{\varepsilon}_i)\) is zero. The map \(\lambda \mapsto \lambda |\mathfrak{h}\) from \(\hat{\mathfrak{h}}_0^*\) to \(\mathfrak{h}^*\) is bijective. Thus, the root system R of \((\mathfrak{g},\mathfrak{h})\) is of type \(A_l\) (Chap. VI, §4, no. 7). Consequently, \(\mathfrak{g}\) is simple (§3, no. 2, Cor. 1 of Prop. 6). Thus, \(\mathfrak{g}\) is a splittable simple Lie algebra of type \(A_l\) .

Every splitting Cartan subalgebra \(\mathfrak{h}^{\prime}\) of \(\mathfrak{g}\) is a transform of \(\mathfrak{h}\) under an elementary automorphism ( \(\S3\) , no. 3, Cor. of Prop. 10). Since \(\operatorname{Aut}_{e}\mathfrak{g}\) is the set of automorphisms \(x \mapsto sxs^{-1}\) of \(\mathfrak{g}\) with \(s \in \mathbf{SL}(V)\) (Chap. VII, \(\S3\) , no. 1, Remark 2; cf.~also (VII)), there exists a basis \(\beta\) of V such that \(\mathfrak{h}^{\prime}\) is the set \(\mathfrak{h}_{\beta}\) of elements of \(\mathfrak{g}\) whose matrix with respect to the basis \(\beta\) is diagonal. Since \(\mathfrak{h}_{\beta}\) contains an element with distinct eigenvalues, the only vector subspaces of V stable under the elements of \(\mathfrak{h}_{\beta}\) are those generated by a subset of \(\beta\) . It follows that the map \(\beta \mapsto \mathfrak{h}_{\beta}\) induces by passage to the quotient a bijection from the set of decompositions of V into the direct sum of \(l + 1\) subspaces of dimension 1 to the set of splitting Cartan subalgebras of \(\mathfrak{g}\).

\begin{enumerate}
\def\labelenumi{(\Roman{enumi})}
\setcounter{enumi}{1}
\tightlist
\item
  Let \(\alpha = \varepsilon_{i} - \varepsilon_{j}\) ( \(i \neq j\) ) be a root. We have \(\mathfrak{g}^{\alpha} = kE_{ij}\) . Since
\end{enumerate}

\[
[ E _ {i j}, E _ {j i} ] = E _ {i i} - E _ {j j}
\]

and since \(\alpha(E_{ii} - E_{jj}) = 2\) , we have (§2, no. 2, Th. 1 (ii))

\[
H _ {\alpha} = E _ {i i} - E _ {j j}.
\]

\begin{enumerate}
\def\labelenumi{(\Roman{enumi})}
\setcounter{enumi}{2}
\tightlist
\item
  Put \(\alpha_{1} = \varepsilon_{1} - \varepsilon_{2},\alpha_{2} = \varepsilon_{2} - \varepsilon_{3},\ldots ,\alpha_{l} = \varepsilon_{l} - \varepsilon_{l + 1}\) . By Chap. VI, §4, no. 7.I, \((\alpha_{1},\dots,\alpha_{l})\) is a basis B of R; the positive roots relative to B are the \(\varepsilon_i - \varepsilon_j\) for \(i < j\) . The corresponding Borel subalgebra \(\mathfrak{b}\) is the set of upper triangular matrices of trace zero.
\end{enumerate}

A flag in V is a set of vector subspaces of V, distinct from \(\{0\}\) and V, totally ordered by inclusion. Order the set of flags of V by inclusion. The maximal flags are the sets \(\{W_{1},\ldots,W_{l}\}\) , where \(W_{i}\) is an i-dimensional vector subspace and

\[
\mathrm{W} _ {1} \subset \dots \subset \mathrm{W} _ {l}.
\]

For example, if \(V_{i}\) denotes the subspace of V generated by \(e_{1},\ldots,e_{i}\) , then \(\{V_{1},\ldots,V_{l}\}\) is a maximal flag.

It is immediate that \(\mathfrak{b}\) is the set of elements of \(\mathfrak{g}\) leaving stable the elements of the maximal flag \(\{V_{1},\ldots,V_{l}\}\) . Conversely, since \(\mathfrak{b}\) contains \(\mathfrak{h}\) and the matrices \(E_{ij}\) for i\textless j, we see that the \(V_{i}\) are the only non-trivial vector subspaces stable under \(\mathfrak{b}\).

Now let \(\delta\) be a maximal flag in V. It follows from the preceding that the set \(\mathfrak{b}_{\delta}\) of elements of \(\mathfrak{g}\) leaving stable all the elements of \(\delta\) is a Borel subalgebra of \(\mathfrak{g}\). Since every Borel subalgebra of \(\mathfrak{g}\) is a transform of \(\mathfrak{b}\) under an elementary automorphism, we see that the map \(\delta \mapsto \mathfrak{b}_{\delta}\) is a bijection from the set of maximal flags to the set of Borel subalgebras of \(\mathfrak{g}\).

Let \(\beta\) be a basis of V. By (I) and the preceding, the Borel subalgebras containing \(\mathfrak{h}_{\beta}\) are those corresponding to the maximal flags each of whose elements is generated by a subset of \(\beta\) . These flags correspond bijectively to the total orders on \(\beta\) in the following way: to a total order \(\omega\) on \(\beta\) is associated the flag \(\{\mathrm{W}_1,\dots ,\mathrm{W}_l\}\) , where \(\mathrm{W}_i\) is the vector subspace generated by the first \(i\) elements of \(\beta\) for the order \(\omega\) . Since there are \((l + 1)!\) total orders on \(\beta\) , we recover the fact that there exist \((l + 1)!\) Borel subalgebras of \((\mathfrak{sl}(\mathrm{V}),\mathfrak{h}_{\beta})\) ( \(\S 3\) , no. 3, Remark).

Let \(\gamma\) be a flag in V. Since \(\gamma\) is contained in a maximal flag, the set \(\mathfrak{p}_{\gamma}\) of elements of \(\mathfrak{g}\) leaving stable the elements of \(\gamma\) is a parabolic subalgebra of \(\mathfrak{g}\). We show that the only non-trivial vector subspaces stable under \(\mathfrak{p}_{\gamma}\) are the elements of \(\gamma\) . For this, we can assume that \(\gamma = \{V_{i_1}, \ldots, V_{i_q}\}\) with \(1 \leq i_1 < \cdots < i_q \leq l\) . Put \(i_0 = 0, i_{q+1} = l + 1\) . The non-empty intervals

\[
\left[ i _ {0} + 1, i _ {1} \right], \left[ i _ {1} + 1, i _ {2} \right], \dots , \left[ i _ {q} + 1, i _ {q + 1} \right]
\]

form a partition of \(\{1, \ldots, l + 1\}\) , so that any square matrix of order \(l + 1\) can be written as a block matrix \((X_{ab})_{1 \leq a, b \leq q + 1}\) . The algebra \(\mathfrak{p}_{\gamma}\) is then the set \(\mathfrak{p}_{i_1,\dots,i_q}\) of elements \((X_{ab})_{1\leq a,b\leq q + 1}\) of \(\mathfrak{sl}(l + 1,k)\) such that \(X_{ab} = 0\) for \(a > b\) . Since \(\mathfrak{p}_{i_1,\dots,i_q}\supset \mathfrak{b}\) , a non-trivial vector subspace stable under \(\mathfrak{p}_{i_1,\dots,i_q}\) is one of the \(\mathrm{V}_i\) ; if \(i_k < i < i_{k + 1}\) , the algebra \(\mathfrak{p}_{i_1,\dots,i_q}\) contains \(E_{i_{k + 1},i}\) and \(\mathrm{V}_i\) is not stable, hence our assertion.

Consequently, the \(2^{l}\) flags contained in the maximal flag \(\{V_{1},\ldots,V_{l}\}\) give rise to \(2^{l}\) distinct parabolic subalgebras containing \(\mathfrak{b}\); since there are exactly \(2^{l}\) parabolic subalgebras containing \(\mathfrak{b}\) ( \(\S3\) , no. 4, Remark), it follows that the map \(\gamma\mapsto \mathfrak{p}_{\gamma}\) is a bijection from the set of flags of V to the set of parabolic subalgebras of \(\mathfrak{g}\). Moreover, \(\mathfrak{p}_{\gamma}\supset \mathfrak{p}_{\gamma'}\) if and only if \(\gamma\subset\gamma'\) .

Recall the parabolic subalgebra \(\mathfrak{p} = \mathfrak{p}_{i_1, \ldots, i_q}\) ( \(1 \leq i_1 < \cdots < i_q \leq l\) ). Let \(\mathfrak{s}\) (resp. \(\mathfrak{n}\)) be the set of \((X_{ab})_{1 \leq a, b \leq q+1}\) in \(\mathfrak{sl}(l+1, k)\) such that \(X_{ab} = 0\) for \(a \neq b\) (resp. \(a \geq b\) ). In view of Prop. 13 of §3, no. 4, we have \(\mathfrak{p} = \mathfrak{s} \oplus \mathfrak{n}\) , the subalgebra \(\mathfrak{s}\) is reductive in \(\mathfrak{g}\) and \(\mathfrak{n}\) is both the largest nilpotent ideal and the nilpotent radical of \(\mathfrak{p}\).

\begin{enumerate}
\def\labelenumi{(\Roman{enumi})}
\setcounter{enumi}{3}
\tightlist
\item
  For \(r = 1, 2, \ldots, l\) , let \(\varpi_{r} = \varepsilon_{1} + \cdots + \varepsilon_{r}\) . We have \(\varpi_{i}(H_{\alpha_{j}}) = \delta_{ij}\) , so \(\varpi_{r}\) is the fundamental weight corresponding to \(\alpha_{r}\) .
\end{enumerate}

Let \(\sigma\) be the identity representation of \(\mathfrak{g}\) on V. The exterior power \(\bigwedge^{r}\sigma\) of \(\sigma\) is a representation on \(\mathrm{E} = \bigwedge^{r}(\mathrm{V})\) . Let \((e_1,\ldots ,e_{l + 1})\) be the chosen basis of V. The \(e_{i_1}\wedge \dots \wedge e_{i_r}\) , where \(i_1 < \dots < i_r\) , form a basis of E. If \(h\in \mathfrak{h}\) ,

\[
(\bigwedge^ {r} \sigma) (h). e _ {i _ {1}} \wedge \dots \wedge e _ {i _ {r}} = (\varepsilon_ {i _ {1}} + \dots + \varepsilon_ {i _ {r}}) (h) e _ {i _ {1}} \wedge \dots \wedge e _ {i _ {r}}.
\]

Thus, every weight is of multiplicity 1, \(\varpi_{r}\) is a weight of \(\bigwedge^{r}\sigma\) , and every other weight is of the form \(\varpi_{r}-\mu\) , where \(\mu\) is a positive radical weight. Consequently, \(\varpi_{r}\) is the highest weight of \(\bigwedge^{r}\sigma\) , and \(e_{1}\wedge\cdots\wedge e_{r}\) is a primitive element. By Chap. VI, §4, no. 7.IX, the Weyl group can be identified with the symmetric group of

\[
\{\varepsilon_ {1}, \dots , \varepsilon_ {l + 1} \}.
\]

The orbit of \(\varpi_{r}\) under the Weyl group thus contains all the \(\varepsilon_{i_{1}} + \cdots + \varepsilon_{i_{r}}\) with \(i_{1} < \cdots < i_{r}\) . The simple submodule generated by the primitive element \(e_{1} \wedge \cdots \wedge e_{r}\) thus admits all the \(\varepsilon_{i_{1}} + \cdots + \varepsilon_{i_{r}}\) as weights and consequently is equal to E. Thus, \(\bigwedge^{r} \sigma\) is irreducible with highest weight \(\varpi_{r}\) .

Thus, the representations \(\bigwedge^{r}\sigma\) ( \(1 \leq r \leq l\) ) are the fundamental representations. We have \(\dim(\bigwedge^{r}\sigma) = \binom{l+1}{r}\) .

\begin{enumerate}
\def\labelenumi{(\Alph{enumi})}
\setcounter{enumi}{21}
\tightlist
\item
  We have \(w_0(\alpha_1) = -\alpha_l, w_0(\alpha_2) = -\alpha_{l-1}, \ldots\) (Chap. VI, §4, no. 7, XI), so
\end{enumerate}

\[
- w _ {0} (\varpi_ {1}) = \varpi_ {l}, - w _ {0} (\varpi_ {2}) = \varpi_ {l - 1}, \ldots .
\]

Let

\[
\omega = n _ {1} \varpi_ {1} + \dots + n _ {l} \varpi_ {l} \quad (n _ {1}, \ldots , n _ {l} \in \mathbf {N})
\]

be a dominant weight. Then, the simple representation with highest weight \(\omega\) is orthogonal or symplectic if and only if

\[
n _ {1} = n _ {l}, \quad n _ {2} = n _ {l - 1}, \ldots
\]

(§7, no. 5, Prop. 12). In particular, if \(l\) is even, none of the fundamental representations of \(\mathfrak{sl}(l + 1,k)\) is orthogonal or symplectic. If \(l\) is odd, the representation \(\bigwedge^i\sigma\) for \(i\neq (l + 1) / 2\) is neither orthogonal nor symplectic; by Chap. VI, §4, no. 7.VI, the sum of the coordinates of \(\varpi_{(l + 1) / 2}\) with respect to \((\alpha_{1},\ldots ,\alpha_{l})\) is

\[
\begin{array}{r l} \frac {1}{l + 1} \left[ \frac {l + 1}{2} \left(1 + 2 + \dots + \frac {l - 1}{2}\right) + \frac {l + 1}{2} \left(1 + 2 + \dots + \frac {l + 1}{2}\right) \right] & = 1 + 2 + \dots + \frac {l - 1}{2} + \frac {l + 1}{4} \end{array}
\]

so \(\bigwedge^{(l + 1) / 2}\sigma\) is orthogonal if \(l\equiv -1\) (mod. 4) and symplectic if \(l\equiv 1\) (mod. 4) (§7, no. 5, Prop. 12). This last result can be made more precise as follows. Choose a non-zero element \(e\) in \(\bigwedge^{l + 1}(\mathrm{V})\) . The multiplication in the exterior algebra \(\mathrm{V}\) defines a bilinear map from

\[
\bigwedge^ {(l + 1) / 2} (\mathrm{V}) \times \bigwedge^ {(l + 1) / 2} (\mathrm{V})
\]

to \(\bigwedge^{l + 1}(\mathrm{V})\) , which can be written \((u,v)\mapsto \varPhi (u,v)e\) , where \(\varPhi\) is a bilinear form on \(\bigwedge^{(l + 1) / 2}(\mathrm{V})\) . It is immediately verified that \(\varPhi\) is non-zero, invariant under \(\mathfrak{g}\) (and hence non-degenerate), symmetric if \((l + 1) / 2\) is even, and alternating if \((l + 1) / 2\) is odd.

\begin{enumerate}
\def\labelenumi{(\Roman{enumi})}
\setcounter{enumi}{5}
\tightlist
\item
  For all \(x \in \mathfrak{g}\) , the characteristic polynomial of \(\sigma(x) = x\) can be written
\end{enumerate}

\[
\mathrm{T} ^ {l + 1} + f _ {2} (x) \mathrm{T} ^ {l - 1} + f _ {3} (x) \mathrm{T} ^ {l - 2} + \dots + f _ {l + 1} (x)
\]

where \(f_{2}, \ldots, f_{l+1}\) are polynomial functions invariant under \(\mathfrak{g}\) (§8, no. 3, Lemma 2).

If \(x = \xi_{1}E_{11} + \cdots + \xi_{l+1}E_{l+1 l+1} \in \mathfrak{h}\) , the \(f_{i}(x)\) are, up to sign, the elementary symmetric functions of \(\xi_{1}, \ldots, \xi_{l+1}\) of degree \(2, \ldots, l + 1\) . Thus, by Chap. VI, §4, no. 7.IX, the \(f_{i}|_{\mathfrak{h}}\) generate the algebra of elements of \(\mathbf{S}(\mathfrak{h}^{*})\) invariant under the Weyl group, and are algebraically independent. Hence (§8, no. 3, Prop. 3) \(f_{2}, f_{3}, \ldots, f_{l+1}\) generate the algebra of polynomial functions invariant under \(\mathfrak{g}\), and are algebraically independent.

\begin{enumerate}
\def\labelenumi{(\Roman{enumi})}
\setcounter{enumi}{6}
\tightlist
\item
  For all \(g \in \mathbf{GL}(l + 1, k)\) , let \(\varphi_k(g) = \varphi(g)\) be the automorphism \(x \mapsto gxg^{-1}\) of \(\mathfrak{g}\) . Then \(\varphi\) is a homomorphism from \(\mathbf{GL}(l + 1, k)\) to \(\operatorname{Aut}(\mathfrak{g})\) . We have
\end{enumerate}

\[
\varphi (\mathbf {S L} (l + 1, k)) = \operatorname{Aut} _ {e} (\mathfrak {g})
\]

(Chap. VII, §3, no. 1, Remark 2). Let \(\bar{k}\) be an algebraic closure of \(k\) . We have

\[
\mathbf {G L} (l + 1, \bar {k}) = \bar {k} ^ {*}. \mathbf {S L} (l + 1, \bar {k}),
\]

so \(\varphi_{\bar{k}}(\mathbf{GL}(l + 1,\bar{k})) = \varphi_{\bar{k}}(\mathbf{SL}(l + 1,\bar{k})) = \mathrm{Aut}_e(\mathfrak{g}\otimes_k\bar{k})\) ; it follows that \(\varphi (\mathbf{GL}(l + 1,k))\subset \mathrm{Aut}_0(\mathfrak{g})\) . On the other hand, \(\mathrm{Aut}_0(\mathfrak{g})\subset \varphi (\mathbf{GL}(l + 1,k))\) , by Prop. 2 of §7, no. 1, applied to the identity representation of \(\mathfrak{g}\) . Hence,

\[
\operatorname{Aut} _ {0} (\mathfrak {g}) = \varphi (\mathbf {G L} (l + 1, k)).
\]

The kernel of \(\varphi\) is the set of elements of \(\mathbf{GL}(l+1,k)\) that commute with every matrix of order \(l+1\) , that is the set \(k^{*}\) of invertible scalar matrices. Thus, \(\operatorname{Aut}_{0}(\mathfrak{g})\) can be identified with the group \(\mathbf{GL}(l+1,k)/k^{*}=\mathbf{PGL}(l+1,k)\) . The kernel of \(\varphi'=\varphi|\mathbf{SL}(l+1,k)\) is \(\mu_{l+1}(k)\) , where \(\mu_{l+1}(k)\) denotes the set of \((l+1)\) th roots of unity in k. Thus, \(\operatorname{Aut}_{e}(\mathfrak{g})\) can be identified with the group \(\mathbf{SL}(l+1,k)/\mu_{l+1}(k)=\mathbf{PSL}(l+1,k)\) . On the other hand, we have the exact sequence

\[
1 \longrightarrow \mathbf {S L} (l + 1, k) \longrightarrow \mathbf {G L} (l + 1, k) \stackrel {\det} {\longrightarrow} k ^ {*} \longrightarrow 1
\]

and the image of \(k^*\) under det is \(k^{*l + 1}\) . It follows that there are canonical isomorphisms

\[
\begin{array}{c} \operatorname{Aut} _ {0} (\mathfrak {g}) / \operatorname{Aut} _ {e} (\mathfrak {g}) \longrightarrow \mathbf {P G L} (l + 1, k) / \mathbf {P S L} (l + 1, k) \\ \qquad \qquad \qquad \qquad \qquad \longrightarrow \mathbf {G L} (l + 1, k) / k ^ {*}. \mathbf {S L} (l + 1, k) \longrightarrow k ^ {*} / k ^ {*   l + 1}. \end{array}
\]

If \(k = \mathbf{R}\) , we see that \(\mathrm{Aut}_0(\mathfrak{g}) = \mathrm{Aut}_e(\mathfrak{g})\) if \(l + 1\) is odd, and that \(\mathrm{Aut}_0(\mathfrak{g}) / \mathrm{Aut}_e(\mathfrak{g})\) is isomorphic to \(\mathbf{Z}/2\mathbf{Z}\) if \(l + 1\) is even.

With the notations of §5, \(f(\mathrm{T}_{\mathbf{Q}})\) is the set of automorphisms of \(\mathfrak{g}\) that induce the identity on \(\mathfrak{h}\), and hence is equal to \(\varphi(\mathrm{D})\) , where D is the set of diagonal elements of \(\mathbf{GL}(l+1,k)\) (§5, Prop. 4). Let \(D'\) be the set of diagonal elements of \(\mathbf{SL}(l+1,k)\) . By Prop. 3 of §5, and the determination of \(\operatorname{Aut}_{e}(\mathfrak{g})\) , we have \(f(q(\mathrm{T}_{\mathrm{P}})) \subset \varphi(\mathrm{D}')\) . We show that \(f(q(\mathrm{T}_{\mathrm{P}})) = \varphi(\mathrm{D}')\) . Let

\[
d = \left( \begin{array}{c c c} \lambda_ {1} & & 0 \\ & \ddots & \\ 0 & & \lambda_ {l + 1} \end{array} \right)
\]

be an element of \(D'\) . There exists a \(\zeta \in \operatorname{Hom}(\mathbf{Q}(\mathbf{R}), k^{*}) = \mathrm{T}_{\mathbf{Q}}\) such that \(\zeta(\varepsilon_{i} - \varepsilon_{j}) = \lambda_{i}\lambda_{j}^{-1}\) for all i and j. It is easy to verify that \(f(\zeta) = \varphi(d)\) . By Chap. VI, §4, no. 7.VIII, \(\mathrm{P}(\mathrm{R})\) is generated by \(\mathrm{Q}(\mathrm{R})\) and the element \(\varepsilon = \varepsilon_{1}\) , whose image in \(\mathrm{P}(\mathrm{R})/\mathrm{Q}(\mathrm{R})\) is of order \(l + 1\) ; but

\[
\begin{array}{c} \zeta ((l + 1) \varepsilon) = \zeta ((\varepsilon_ {1} - \varepsilon_ {2}) + (\varepsilon_ {1} - \varepsilon_ {3}) + \dots + (\varepsilon_ {1} - \varepsilon_ {l + 1})) \\ = \lambda_ {1} ^ {l} \lambda_ {2} ^ {- 1} \lambda_ {3} ^ {- 1} \ldots \lambda_ {l + 1} ^ {- 1} = \lambda_ {1} ^ {l + 1} \end{array}
\]

so \(\zeta\) extends to a homomorphism from \(\mathrm{P}(\mathrm{R})\) to \(k^{*}\) . This proves that \(\zeta \in q(\mathrm{T}_{\mathrm{P}})\) , so \(\varphi(d) \in f(q(\mathrm{T}_{\mathrm{P}}))\) .

Recall (§5, no. 3, Cor. 2 of Prop. 5) that \(\mathrm{Aut}(\mathfrak{g}) = \mathrm{Aut}_0(\mathfrak{g})\) for \(l = 1\) , and that \(\mathrm{Aut}(\mathfrak{g}) / \mathrm{Aut}_0(\mathfrak{g})\) is isomorphic to \(\mathbf{Z}/2\mathbf{Z}\) for \(l \geq 2\) . The map \(\theta : x \mapsto -^t x\) is an automorphism of \(\mathfrak{sl}(l + 1, k)\) and \(a_0 = \theta|\mathfrak{h} \notin \mathrm{W}\) if \(l \geq 2\) (Chap. VI, §4, no. 7.XI), so the class of \(a_0\) in \(\mathrm{Aut}(\mathfrak{g}) / \mathrm{Aut}_0(\mathfrak{g})\) is the non-trivial element of this group (§5, no. 2, Prop. 4).

\begin{enumerate}
\def\labelenumi{(\Roman{enumi})}
\setcounter{enumi}{7}
\tightlist
\item
  The restriction to \(\mathfrak{h}\) of the Killing form is
\end{enumerate}

\[
\begin{array}{l} \Phi (\xi_ {1} E _ {1 1} + \dots + \xi_ {l + 1} E _ {l + 1, l + 1}, \xi_ {1} ^ {\prime} E _ {1 1} + \dots + \xi_ {l + 1} ^ {\prime} E _ {l + 1, l + 1}) \\ = \sum_ {i \neq j} (\xi_ {i} - \xi_ {j}) (\xi_ {i} ^ {\prime} - \xi_ {j} ^ {\prime}) = \sum_ {i, j} (\xi_ {i} - \xi_ {j}) (\xi_ {i} ^ {\prime} - \xi_ {j} ^ {\prime}) \\ = (l + 1) \sum_ {i} \xi_ {i} \xi_ {i} ^ {\prime} + (l + 1) \sum_ {j} \xi_ {j} \xi_ {j} ^ {\prime} - 2 \left(\sum_ {i} \xi_ {i}\right) \left(\sum_ {j} \xi_ {j} ^ {\prime}\right) \\ = 2 (l + 1) \sum_ {i} \xi_ {i} \xi_ {i} ^ {\prime}. \end{array}
\]

\begin{enumerate}
\def\labelenumi{(\Roman{enumi})}
\setcounter{enumi}{8}
\tightlist
\item
  For \(1 \leq i < j \leq l + 1\) , put
\end{enumerate}

\[
X _ {\varepsilon_ {i} - \varepsilon_ {j}} = E _ {i j} X _ {\varepsilon_ {j} - \varepsilon_ {i}} = - E _ {j i}.
\]

Then, for all \(\alpha \in \mathrm{R}\) , we have \([X_{\alpha}, X_{-\alpha}] = -H_{\alpha}\) and \(\theta(X_{\alpha}) = X_{-\alpha}\) (where \(\theta\) is the automorphism \(x \mapsto -^{t} x\) introduced in (VII)). Consequently, \((X_{\alpha})_{\alpha \in \mathrm{R}}\) is a Chevalley system in \((\mathfrak{g}, \mathfrak{h})\) .

Take \(k = \mathbf{Q}\) . The permissible lattices in \(\mathfrak{h}\) (§12, no. 6, Def. 1) are those lying between the \(\mathbf{Z}\) -module \(\mathrm{Q}(\mathrm{R}^{\vee})\) generated by the \(E_{ii} - E_{i+1,i+1}\) , that is, consisting of the diagonal matrices belonging to \(\mathfrak{sl}(l+1,\mathbf{Z})\) , and the \(\mathbf{Z}\) -module \(\mathrm{P}(\mathrm{R}^{\vee})\) generated by \(\mathrm{Q}(\mathrm{R}^{\vee})\) and \(E_{11} - (l+1)^{-1}\sum E_{ii}\) (Chap. VI, §4, no. 7.VIII), that is, consisting of the diagonal matrices of trace zero of the form \(x + (l+1)^{-1}a.1\) , where \(x\) has integer entries and \(a\in\mathbf{Z}\) . It follows that \(\mathfrak{sl}(l+1,\mathbf{Z})\) is the Chevalley order in \((\mathfrak{g},\mathfrak{h})\) associated to the permissible lattice \(\mathrm{Q}(\mathrm{R}^{\vee})\) and the Chevalley system \((X_{\alpha})\) . It is easy to verify that \(\bigwedge^{r}\mathbf{Z}^{l+1}\) is an admissible lattice in \(\bigwedge^{r}\mathbf{Q}^{l+1}\) relative to \(\mathfrak{sl}(l+1,\mathbf{Z})\) (§12, no. 8, Def. 3).

On the other hand, \(\mathfrak{gl}(l+1,\mathbf{Z})\) is a Chevalley order in the split reductive algebra \(\mathfrak{gl}(l+1,\mathbf{Q})\) ; its projection onto \(\mathfrak{sl}(l+1,\mathbf{Q})\) parallel to the centre \(\mathbf{Q}.1\) of \(\mathfrak{gl}(l+1,\mathbf{Q})\) is the Chevalley order in \((\mathfrak{g},\mathfrak{h})\) defined by the permissible lattice \(\mathrm{P}(\mathrm{R}^{\vee})\) in \(\mathfrak{h}\) and the Chevalley system \((X_{\alpha})\) . We remark that \(\mathfrak{gl}(l+1,\mathbf{Z})\) is not the direct sum of its intersections with \(\mathfrak{sl}(l+1,\mathbf{Q})\) and the centre of \(\mathfrak{gl}(l+1,\mathbf{Q})\) .
% semantic-fragments: legacy-input-seam
\relax{}
